% Pascal's theorem: the opposite sides g_i = P_i P_{i+1} and g_{i+3} of a
% hexagon inscribed in a conic meet in three collinear points.
\begin{tikzpicture}[x=0.6cm, y=0.6cm]
  \clip (0.07,-5.24) rectangle (14.04,2.69);
  \draw[curve component, rotate around={-177.03:(7.15,-1.27)}]
    (7.15,-1.27) ellipse [x radius=4.91, y radius=2.28];
  \foreach \i/\x/\y in {1/2.26/-1.64, 2/7.92/1.02, 3/9.8/-3.08, 4/4.86/0.66,
      5/6.72/-3.56, 6/10.58/0.5}
    \coordinate (P\i) at (\x,\y);
  \foreach \a/\b in {1/2, 2/3, 3/4, 4/5, 5/6, 6/1}
    \draw[name path global=g\a] ($(P\a)!-4!(P\b)$) -- ($(P\b)!-4!(P\a)$);
  \foreach \i/\x/\y in {1/1/-2.825, 2/6.8/2.1, 3/12.05/-4, 4/5.05/1.85,
      5/5.1/-4.5, 6/13/0.6}
    \node at (\x,\y) {$g_{\i}$};
  \path[name intersections={of=g1 and g4, by=Q1}];
  \path[name intersections={of=g2 and g5, by=Q2}];
  \path[name intersections={of=g3 and g6, by=Q3}];
  \draw[dzg accent, dashed] ($(Q1)!-8!(Q2)$) -- ($(Q2)!-8!(Q1)$);
  \foreach \i/\pos in {1/left, 2/above left, 3/below right, 4/above left,
      5/below, 6/above right}
    \node[curve point, label={[fill=white, inner sep=1pt]\pos:$P_{\i}$}] at (P\i) {};
  \foreach \i in {1,2,3} \node[constructed point] at (Q\i) {};
\end{tikzpicture}
